\section{Closure, one phase and recurring wave}
\label{sec:closure-phase-wave}

\subsection{Completed retained return}
\label{subsec:monon-beat}
\label{subsec:monon-cycle}
\label{app:wave}
\label{app:wave:scope}
A closure trace retains its start even as the current crown advances:
\begin{equation}
\Gamma_C=(B,r,i,a_0,(e_0,\ldots,e_n),M_{\varepsilon,m},\mathcal W_C).
\label{eq:closure-trace}
\end{equation}
Here \(r\) is scale, \(i\) is lineage, the events are strictly committed in
order, and \(\mathcal W_C\) binds the hinge and return to the retained start.
All traversals satisfy \eqref{eq:atomic-h1}. The closure certificate is
\begin{equation}
\operatorname{Cert}(C)=
(B,r,i,\operatorname{id}(e_0),\operatorname{id}(e_n),
\operatorname{digest}\Gamma_C,\mathcal W_C,\mathrm{closed}),
\quad \operatorname{end}(C)\sim_{B,i,\mathcal W_C}\operatorname{start}(C).
\label{eq:closure-certificate}
\end{equation}
The equivalence preserves the bound start and ancestry. Equal unbound
endpoints alone supply no certificate. A monon is the primitive completed
closure of \eqref{eq:hinge-grammar}, written
\begin{equation}
\mathrm{Monon}(x)=\operatorname{cycle}_{\mathsf H_3}(x\succ x=x).
\label{eq:monon-cycle}
\end{equation}
Completion contributes one phase:
\begin{equation}
\Delta\Phi(C)=1_C.
\label{eq:closure-phase}
\end{equation}
Additional dwell does not multiply this count. The two reflected legs are
parts of the same completed phase. An isotropic or anisotropic model can
support a certified closure; its oriented ancestral trace remains retained.

\paragraph{Symmetry.}
\begin{equation}
\operatorname{Close}(C)=1\;\not\Rightarrow\;\operatorname{Isotropic}(C).
\label{eq:closure-not-isotropy}
\end{equation}
Isotropy refers to a named support, model and relabelling action. Cancellation
of aggregate values preserves the contributing oriented histories. Neither
symmetry nor anisotropy changes \eqref{eq:closure-phase}.

\subsection{Advance and closure memory}
\label{subsec:crown}
An admitted successor becomes the next crown through a committed advance:
\begin{equation}
\operatorname{Adv}_k(x_k^0)=x_{k+1}^-,\qquad
\operatorname{Adv}_k(h_k^+)=x_{k+1}^0,\qquad h_k^+\in\mathcal H_k^+.
\label{eq:crown-advance}
\end{equation}
Advance carries the selected member into the next crown. A later
lineage-bound return can re-crown the retained start:
\begin{equation}
\operatorname{Crown}(\mathsf H_3^{\mathrm{role}}(k+L))
\sim_{B,i,\mathcal W_C}\operatorname{Crown}(\mathsf H_3^{\mathrm{role}}(k)),\qquad L>0.
\label{eq:closure-memory}
\end{equation}
Open advance is a continuation; the further return witness establishes closure.
Solitonic recurrence names preservation of certified relational form through
the stated advance and reclosure, with the ancestry and scope map retained.

\subsection{Return identity, beat and cadence}
\label{app:wave:cycle-wave-identity}
Give each completed return immutable identity
\begin{equation}
\operatorname{id}_{\mathrm{return}}(C)=
(B,r,i,\operatorname{id}(e_0),\operatorname{id}(e_n),
\operatorname{digest}\Gamma_C).
\label{eq:return-identity}
\end{equation}
Certificate wrapping, repeated observation of the same closed state, and
replay of that return retain this identity. A distinct compatible certified
reclosure after completion contributes one beat. With ordered completion
indices \(k_0<k_1<\cdots<k_m\), define
\begin{equation}
\mathrm{bip}_i(\omega)=
\#\{C_{i,j}:j\ge1,\ k_j\in\omega,\ \operatorname{id}_{\mathrm{return}}(C_{i,j})
\text{ is distinct}\},\qquad
\operatorname{Sep}_{i,j}=k_{j+1}-k_j>0.
\label{eq:beat-reclosure}
\end{equation}
The initial closure supplies the origin of this reclosure count. At an integer
cut \(k\), the cumulative accessor counts distinct committed return identities:
\begin{equation}
\begin{aligned}
N_i^C(k)&=\#\{\operatorname{id}_{\mathrm{return}}(C_{i,j}):k_j\le k\},\\
\Phi_i(k)&=N_i^C(k)1_C,\qquad
\mathrm{bip}^{\le}_i(k)=\max(N_i^C(k)-1,0).
\end{aligned}
\label{eq:cumulative-closure-counts}
\end{equation}
Thus \(N\ge1\) completed closures give \(N\) phase units and \(N-1\)
reclosure beats from the lineage origin. For an arbitrary native window
\(\omega=(u,v]\), \(u,v\in\mathbb Z\), \(u<v\), the same distinct-return count gives
\begin{equation}
\Phi_i(v)-\Phi_i(u)=
\bigl(\mathrm{bip}_i((u,v])+\mathbf1_{\{u<k_0\le v\}}\bigr)1_C.
\label{eq:closure-phase-window}
\end{equation}
The origin contributes the indicator; later windows do not subtract another
initial closure. The exact ordered separations form cadence \(\operatorname{Cad}_i=(\operatorname{Sep}_{i,j})_j\).
It may be nonuniform; RD support and harmonic display period are separate
quantities. Recurrence is not inferred from a repeated value alone.

\paragraph{Recurring wave.}
For a compatible certified closure/current lineage with distinct reclosure,
fix immutable wave identity and its growing admitted history:
\begin{equation}
W^{\mathrm{id}}_{B,i}=(B,r,i,K_i^{\mathrm{id}},
\operatorname{id}_{\mathrm{return}}(C_{i,0})),\qquad
W_{B,i}|_k=(W^{\mathrm{id}}_{B,i},(C_{i,j})_{k_j\le k},\operatorname{Cad}_i|_k).
\label{eq:wave-identity}
\end{equation}
Compatibility retains the origin's ancestral identity at fixed scope and
representation; changed representations require their declared identity-preserving
scope map. At least two compatible distinct completed closures precede recurrence. A single closure is one phase; recurring retained
closure/current with cadence is wave. Later field organization resolves these
already admitted identities.

\paragraph{Native running and standing classifications.}
A running-wave classification retains an oriented recurring current lineage;
recurring oriented Duon current is the inherited running-Duon construction.
A standing organization binds compatible oppositely oriented recurrent
lineages under a declared coupling and balance condition. Zero aggregate
current by itself supplies neither recurrence nor that coupling witness.
Both classifications retain their individual return identities and histories;
the character and running/standing coordinates in \eqref{eq:fourier-transform}
and \eqref{eq:fourier-running-standing} represent these admitted organizations.\label{eq:native-running-standing}

\subsection{Reflected-leg comparison}
\label{subsec:reflected-leg-phase}
\paragraph{Setup.}
Declare a uniform atomic-event support comparison, with formal quantum
\(h\) in the ordered additive domain \(\mathcal L=\mathbb N_0 h,\ h>0\).
A target closure and a primitive reference are already certified at the same
scope and scale. Their reflection witnesses bind the outbound events, committed
hinge/turnaround, returning events and retained target start. Both legs are
nonempty; an arbitrary cut of a word is not that witness. Count each atomic
traversal once, excluding hinge dwell and replay:
\begin{equation}
\delta\phi^{\rightarrow}=n_{\rightarrow}h,\qquad
\delta\phi^{\leftarrow}=n_{\leftarrow}h,\qquad
n_{\rightarrow},n_{\leftarrow}\in\mathbb Z_{\ge1}.
\label{eq:directed-leg-operands}
\end{equation}
The reference is a certified two-edge retained-start return with mandatory
hinge. Its support is \(2h\). Fix the comparison embedding before examining
the target:
\begin{equation}
\iota_{\mathrm{ref}}:\mathbb N_0 1_C\longrightarrow\mathcal L,
\qquad \iota_{\mathrm{ref}}(n1_C)=2nh.
\label{eq:phase-unit-embedding}
\end{equation}
This embedding compares two distinct domains: closure count and declared
event support. Uniform weighting and this reference are model choices.

\paragraph{Derivation.}
Each committed leg has at least one atomic traversal, so addition in
\(\mathcal L\) gives
\begin{equation}
\delta\phi^{\rightarrow}+\delta\phi^{\leftarrow}
=(n_{\rightarrow}+n_{\leftarrow})h
\ge2h=\iota_{\mathrm{ref}}(1_C).
\label{eq:reflected-leg-bound}
\end{equation}
One-phase normalization by itself does not prove this support inequality.
Its premises are the explicit event convention and certified reference.

\paragraph{Specialization.}
Under declared reflected-path symmetry and equal operands,
\begin{equation}
\delta\phi^{\rightarrow}=\delta\phi^{\leftarrow}=\delta\phi
\quad\Longrightarrow\quad
\delta\phi^{\rightarrow}+\delta\phi^{\leftarrow}=2\delta\phi
\ge\iota_{\mathrm{ref}}(1_C).
\label{eq:equal-leg-factorization}
\end{equation}
An asymmetric \(1+3\) trace has support \(4h\) and one closure phase. Its sum
equals neither \(2h\) nor \(6h\), the doubled individual legs. In this uniform
model each doubled leg may separately satisfy the lower bound; equality is
required for factorization of the sum into one common operand. With the
embedding understood, \(2\delta\phi\ge1_C\) denotes this specialization.

\paragraph{Falsification.}
A missing or misbound reflection witness, signed branch cancellation used as
support, incompatible reference scale, nonminimal reference, or unequal-leg
factorization violates the derivation. This comparison supplies no pre-closure
admission selector, duration, stochastic weight or flux.

\paragraph{Provenance.}
Equations \eqref{eq:atomic-h1}, \eqref{eq:hinge-grammar},
\eqref{eq:closure-certificate}, and the declared comparison
\eqref{eq:phase-unit-embedding}. The Manual executes symmetric and asymmetric
lineage-bound cases with separate phase and support counts.

\paragraph{Continuation.}
\label{app:wave:fractal-growth-temporal-direction}
Hinge saturation restricts further same-resolution dwell. A later hinge and
an admitted continuation may follow; only a certified retained-start return
adds closure, and only a distinct compatible reclosure adds recurrence.
\paragraph{Localized recurrence.}
\label{subsec:localized-wavelet}
A native wavelet is a localized segment of an already admitted recurrent wave:
it retains that wave's identity, its adjacent certified return events and their
positive local separation. Localization selects a segment of retained recurrence;
it does not promote an isolated monon or an arbitrary scalar pulse into a wave.
Its later harmonic representation, when used, binds that same finite segment.
\paragraph{Falsification of wave admission.}
\label{app:wave:falsification}
An isolated closure, a replayed certificate, an unbound lineage, or an
unrealized successor horizon cannot satisfy \eqref{eq:wave-identity}.
